\begin{abstract}
Fine-grained recognition of cat and dog breeds is challenging because visually similar classes differ primarily in subtle, localized anatomical cues, while unconstrained real-world images suffer from background clutter, multiple animals, and non-target distractors. We address a 20-class breed recognition task with a dedicated reject class using a lightweight two-stage cascade pipeline. A COCO-pretrained YOLOv6 detector localizes subjects, isolates the primary animal by bounding-box area, masks secondary detections to eliminate context leakage, and extracts a standardized $224\times224$ crop. Species-based routing directs the crop to one of two domain-specific EfficientNetV2-S classifiers trained from scratch for ten cat and ten dog breeds, respectively. Empirical evaluation across multiple augmentation strategies demonstrates that spatial patch-mixing (CropMix) acts as the primary regularizer, pushing dog classification accuracy from a $78.66\%$ unregularized baseline to $85.02\%$ while eliminating validation loss divergence. Furthermore, xAi diagnostics via Grad-CAM confirm that spatial masking and heavy augmentation successfully suppress background shortcut learning, forcing the network to focus on discriminative morphological traits. End-to-end performance benchmarking yields a mean per-image inference latency of $200.63~\text{ms}$, utilizing only $4.01\%$ of the $5.0$-second runtime constraint. In the end we reached an overall accuracy of 92.35\%.
\end{abstract}